% ============================================================
% SECȚIUNEA 5 — Securitate
% Fișiere sursă de referință:
%   backend/AIInterviewCoach.Infrastructure/Services/JwtTokenService.cs
%   backend/AIInterviewCoach.Infrastructure/Services/PasswordHasherService.cs
%   backend/AIInterviewCoach.API/Authentication/AuthCookieDefaults.cs
%   backend/AIInterviewCoach.API/Extensions/ServiceCollectionExtensions.cs (AddJwtAuthentication)
%   backend/AIInterviewCoach.API/Authorization/AuthorizationPolicies.cs
%   backend/AIInterviewCoach.API/RateLimiting/RateLimitingPolicies.cs
%   backend/AIInterviewCoach.Application/Services/ProblemAccessAuthorization.cs
%   frontend/src/app/core/interceptors/auth.interceptor.ts
%   frontend/src/app/core/guards/auth-guard.ts
%   frontend/src/app/core/guards/role-guard.ts
%   frontend/src/app/core/auth/access-policies.ts
% ============================================================

\section{Securitatea sistemului}
\label{sec:securitate}

\subsection{Autentificarea utilizatorilor}

\subsubsection{Hash-ul parolelor}

Parolele sunt procesate înainte de stocare cu \texttt{PasswordHasherService},
care utilizează \textbf{PBKDF2-HMAC-SHA256}:
% Sursă: Infrastructure/Services/PasswordHasherService.cs

\begin{itemize}
  \item Salt aleatoriu de \textbf{16 octeți}, generat cu
        \texttt{RandomNumberGenerator.GetBytes(16)} (CSP, criptografic sigur).
  \item \textbf{10\,000 de iterații} de derivare a cheii.
  \item Output de \textbf{32 de octeți} (256 biți).
  \item Stocare ca un singur șir \texttt{base64(salt).base64(hash)};
        salt-ul este stocat alături de hash, deci fiecare parola are un
        salt unic.
\end{itemize}

Verificarea folosește \texttt{CryptographicOperations.FixedTimeEquals()},
care compară octeți în timp constant, eliminând vulnerabilitatea de tip
\emph{timing attack}.

\subsubsection{Emiterea token-ului JWT}

La autentificare cu succes, \texttt{JwtTokenService.GenerateToken()} emite
un token JWT semnat \textbf{HMAC-SHA256} cu:
% Sursă: Infrastructure/Services/JwtTokenService.cs

\begin{itemize}
  \item Claim \texttt{sub} — GUID-ul utilizatorului.
  \item Claim \texttt{email} — adresa de e-mail.
  \item Claim \texttt{name} — numele complet.
  \item Claim \texttt{role} — rolul (stringul enum: \texttt{Candidate},
        \texttt{Interviewer} sau \texttt{Admin}).
  \item Câmpuri \texttt{iss} (issuer) și \texttt{aud} (audience) din
        configurare (\texttt{Jwt:Issuer}, \texttt{Jwt:Audience}).
  \item Timp de viață: \textbf{7 zile} (\texttt{DateTime.UtcNow.AddDays(7)}).
\end{itemize}

Cheia de semnare este citită obligatoriu din configurare
(\texttt{JwtSigningKeyResolver.GetRequiredSigningKey()}); dacă lipsește,
aplicația nu pornește.

\subsubsection{Transportul token-ului}

Token-ul este plasat de server în cookie-ul HTTP \texttt{aic\_auth} cu
durata de viață de 7 zile, sau poate fi transmis ca antet HTTP standard
\texttt{Authorization: Bearer <token>}.
% Sursă: API/Authentication/AuthCookieDefaults.cs

Middleware-ul \texttt{OnMessageReceived} din configurarea JwtBearer
implementează o logică de prioritate:
% Sursă: ServiceCollectionExtensions.cs:96–118

\begin{enumerate}
  \item Dacă \texttt{context.Token} este deja populat (alt middleware), nu
        face nimic.
  \item Încearcă cookie-ul \texttt{aic\_auth}.
  \item Dacă cookie-ul lipsește, încearcă antetul \texttt{Authorization}.
\end{enumerate}

Validarea token-ului verifică emitent, audiență, semnătură și termenul de
valabilitate (\texttt{ValidateLifetime = true}).

\subsection{Autorizarea bazată pe roluri}

\subsubsection{Politici de autorizare — backend}

ASP.NET Core evaluează politicile de autorizare după autentificare.
Aplicația definește cinci politici înregistrate în
\texttt{AuthorizationPolicies.Register()}:
% Sursă: API/Authorization/AuthorizationPolicies.cs

\begin{itemize}
  \item \texttt{InterviewerWorkspaceAccess} — acceptă rolurile
        \texttt{Interviewer} și \texttt{Admin}.
        Protejează CRUD-ul de interviuri și gestionarea sesiunilor.
  \item \texttt{CandidateWorkspaceAccess} — acceptă rolurile
        \texttt{Candidate} și \texttt{Admin}.
        Protejează submisiile, accesul prin token și finalizarea sesiunii.
  \item \texttt{AdminProblemManagement} — exclusiv \texttt{Admin}.
        Protejează crearea, modificarea și ștergerea problemelor, a
        cazurilor de test și a catalog-ului.
  \item \texttt{AdminAuditLogAccess} — exclusiv \texttt{Admin}.
        Protejează interogarea jurnalelor de audit.
  \item \texttt{AdminObservabilityAccess} — exclusiv \texttt{Admin}.
        Protejează metricele de observabilitate.
\end{itemize}

\subsubsection{Autorizarea la nivel de date — ProblemAccessAuthorization}

Politicile de mai sus operează la nivel de rută. La nivel de date,
\texttt{ProblemAccessAuthorization.EnsureProblemAccessible()} adaugă un
strat suplimentar:
% Sursă: Application/Services/ProblemAccessAuthorization.cs

\begin{itemize}
  \item \textbf{Admin}: acces necondiționat la orice problemă.
  \item \textbf{Candidate}: poate accesa numai problemele cu
        \texttt{IsPublic = true}.
  \item \textbf{Interviewer}: poate accesa orice problemă publică sau
        orice problemă proprie (\texttt{CreatedByUserId == currentUserId}).
\end{itemize}

O metodă separată, \texttt{EnsureProblemOwnership()}, verifică că un
intervievator modifică/șterge doar propriile cazuri de test (sau că
solicitantul este admin).

\subsubsection{Token-urile de acces la interviuri}

Câmpul \texttt{AccessToken} din tabela \texttt{Interviews} este un șir
opac generat la crearea interviului. Endpoint-ul
\texttt{GET /api/interviews/token/\{token\}} este deliberat \emph{neautentificat},
permițând candidaților să vizualizeze detaliile interviului înainte de
conectare. Endpoint-ul
\texttt{POST /api/interviews/token/\{token\}/start} cere autentificare
(politica \texttt{CandidateWorkspaceAccess}) și creează sesiunea.
Rata de acces prin token este limitată la 20 cereri/minut per utilizator
(politica \texttt{InterviewTokenAccess}), prevenind enumerarea token-urilor.

\subsection{Rate limiting}

Limitele de rată sunt implementate cu
\texttt{Microsoft.AspNetCore.RateLimiting} și partiționate per sursă:
% Sursă: API/RateLimiting/RateLimitingPolicies.cs

\begin{itemize}
  \item \textbf{AuthLogin} (fereastră fixă, 5 cereri/minut per IP):
        Protejează \texttt{POST /api/auth/register} și
        \texttt{POST /api/auth/login} împotriva atacurilor de forță brută.
        Partiționare: adresă IP a clientului.
        \emph{Notă}: în mediul de testare de integrare limita este ridicată
        la 50 cereri/minut.

  \item \textbf{InterviewTokenAccess} (fereastră fixă, 20 cereri/minut per
        utilizator):
        Protejează \texttt{GET /token/\{token\}} și
        \texttt{POST /token/\{token\}/start} împotriva enumerării token-urilor.
        Partiționare: ID utilizator (sau IP dacă neautentificat).

  \item \textbf{CandidateSubmissionFlow} (bucket de tokeni, 15
        tokeni/minut per utilizator):
        Protejează \texttt{POST /submissions}, \texttt{GET /submissions/me},
        \texttt{DELETE /submissions/...} și \texttt{POST /problems/\{id\}/hints}.
        Algoritmul bucket de tokeni permite rafale mici (până la 15 cereri
        instant dacă bucket-ul este plin) fără a penaliza utilizarea normală.
        Partiționare: ID utilizator.

  \item \textbf{AdminMutation} (fereastră fixă, 12 cereri/minut per
        utilizator):
        Protejează operațiunile de scriere administrative (creare/modificare/
        ștergere probleme și cazuri de test, înlocuire catalog).
        Partiționare: ID utilizator.
\end{itemize}

La depășirea limitei, API-ul răspunde cu \textbf{HTTP~429} și, dacă
aplicabil, cu antetul \texttt{Retry-After} exprimat în secunde.

\subsection{Securitate frontend}

\subsubsection{Interceptorul de autentificare}

\texttt{authInterceptor} (Angular \texttt{HttpInterceptorFn}) clonează
fiecare cerere HTTP care vizează \texttt{environment.apiBaseUrl} și
adaugă \texttt{withCredentials: true}, asigurând că browser-ul trimite
automat cookie-ul \texttt{aic\_auth} la fiecare cerere API.
Cererile către alte origini nu sunt modificate.
% Sursă: frontend/src/app/core/interceptors/auth.interceptor.ts

\subsubsection{Gardele de rută}

\begin{itemize}
  \item \texttt{authGuard}: verifică \texttt{AuthService.isAuthenticated()};
        dacă utilizatorul nu este autentificat, redirecționează la
        \texttt{/login}.
        % Sursă: frontend/src/app/core/guards/auth-guard.ts

  \item \texttt{roleGuard}: citește politica de acces din datele rutei
        (\texttt{route.data['accessPolicy']}) și evaluează rolul curent
        al utilizatorului prin \texttt{canAccessPolicy()}.
        Dacă rolul nu îndeplinește politica:
        \begin{itemize}
          \item Utilizator neautentificat $\rightarrow$ redirecționare la \texttt{/login}.
          \item Utilizator cu alt rol $\rightarrow$ redirecționare la ruta implicită
                a rolului (\texttt{getDefaultRouteForRole()}):
                Interviewer/Admin $\rightarrow$ \texttt{/interviewer/dashboard},
                Candidate $\rightarrow$ \texttt{/candidate/access}.
        \end{itemize}
        % Sursă: frontend/src/app/core/guards/role-guard.ts

  \item Funcțiile de politică (\texttt{canAccessInterviewerWorkspace},
        \texttt{canAccessCandidateWorkspace}, \texttt{isAdminRole}) oglindesc
        exact politicile de autorizare definite în backend, asigurând
        consistența verificărilor.
        % Sursă: frontend/src/app/core/auth/access-policies.ts
\end{itemize}

\subsection{Considerații de securitate suplimentare}

\paragraph{Rezistența la timing attacks}
Compararea hash-urilor de parolă folosește \texttt{FixedTimeEquals()},
care execută întotdeauna același număr de operații indiferent de lungimea
prefixului comun, eliminând posibilitatea inferenței prin variația
timpului de răspuns.
% Sursă: Infrastructure/Services/PasswordHasherService.cs

\paragraph{Separarea responsabilităților autentificare/autorizare}
Token-urile de acces la interviu (\texttt{AccessToken}) sunt separate
de autentificarea cu parolă: un candidat poate vedea un interviu fără
cont, dar nu poate \emph{începe} sesiunea fără autentificare. Aceasta
permite unui intervievator să trimită linkul de interviu pe e-mail înainte
ca acel candidat să fie înregistrat în sistem.

\paragraph{Jurnalul de audit administrativ}
Toate acțiunile administrative sensibile (creare/modificare/ștergere
probleme, înlocuire catalog) sunt înregistrate în tabela
\texttt{AdminAuditLogs} cu ID-ul, e-mailul și numele complet al
administratorului, tipul acțiunii, entitatea afectată și rezumatul.
Aceasta asigură trasabilitatea modificărilor și suportul pentru
investigații post-incident.
% Sursă: Domain/Entities/AdminAuditLog.cs
